\begin{proposition}[有限周期驱动的同步窗口]
对一个 Bloch 场分量 $u_m\ee^{\ii(k_z+G_m)z}$，在长度 $L$ 的均匀窗口内，轨迹振幅因子为
\begin{equation}
\int_{-L/2}^{L/2}\ee^{\ii\Delta k z}\dd z
=L\sinc(\Delta kL/2),
\qquad\Delta k=k_z+G_m-\omega/v.
\label{eq:feqo-periodic-phase-matching}
\end{equation}
其强度主瓣第一零点为 $|\Delta k|=2\pi/L$，全半高宽为 $5.56623\ldots/L$。同步是允许条件；若模式重叠 $u_m=0$，该通道仍无激发。
\end{proposition}
\begin{proof}
$\Delta k\ne0$ 时直接积分指数得 $2\sin(\Delta kL/2)/\Delta k$；零点处由连续性得到 $L$。强度归一化后为 $\sinc^2(\Delta kL/2)$，其第一半高根 $x_{1/2}=1.391557\ldots$ 给出全宽 $4x_{1/2}/L$。
\end{proof}

\begin{workedexample}[垂直出射与有限角度宽度]
对真空出射光 $k_z=(\omega/c)\cos\theta$，小反冲下的同步式恢复\eqnrefs{eq:smith_purcell_wavelength}。在 $\theta=\pi/2$ 时，周期 $\Lambda=m\lambda_{\rm rad}v/c$。若光波长为 $800\,\mathrm{nm}$、$v=0.70c$、$m=1$，周期为 $560\,\mathrm{nm}$。

固定频率时，中心到第一零点的方向余弦差为 $|\Delta\cos\theta|=\lambda_{\rm rad}/L$；强度全半高宽为 $5.56623\lambda_{\rm rad}/(2\pi L)$。这些是理想均匀窗口的宽度，还需乘模式耦合和收集效率。

无限匀速亚光速电子的同步线位于真空光锥外。周期倒格矢是把它搬入传播区的一种机制；有限路径、加速与边缘也会改变源谱或平移对称性，应由相应边界问题计算。
\end{workedexample}

\begin{figure}[htbp]
\centering
\includegraphics[width=.72\linewidth]{light_cone_synchronism.pdf}
\caption{亚光速电子同步线与真空光锥。无限平移对称几何中，两者除原点外不相交；周期结构可通过倒格动量平移源谱，有限路径及加速则需使用对应的源与边界条件。}
\label{fig:light-cone-synchronism}
\end{figure}
